\section*{Hoofdstuk \ref{ch:b3:extension-strike-slip} --- Extensie- en horizontale-verschuivingstektoniek}

\begin{solution}{exo:b3:extension-strike-slip:1}
De lagen hellen $50 - 35 = \ang{15}$, weg van de breuken; $\beta =
\sin\ang{50}/\sin\ang{35} = 0.766/0.574 = 1.34$.
\end{solution}

\begin{solution}{exo:b3:extension-strike-slip:2}
Dextrale stap naar rechts: bekken. Dextrale stap naar links: rug. Sinistrale
stap naar links: bekken. Sinistrale stap naar rechts: rug.
\end{solution}

\begin{solution}{exo:b3:extension-strike-slip:3}
$105/20 = \qty{5.25}{km/Myr} = \qty{5.25}{mm/yr}$, dicht bij de
\qty{6}{mm/yr}: de snelheid hoeft over \qty{20}{Myr} niet constant te zijn geweest.
\end{solution}

\begin{solution}{exo:b3:extension-strike-slip:4}
\qty{200}{\percent} betekent $\Delta L = 2L_0$: $L_0 = \qty{117.5}{km}$ en $L
= \qty{352.5}{km}$. In het noorden geldt $\Delta L = 0.5L_0$: $L_0 = \qty{470}{km}$
en $L = \qty{705}{km}$.
\end{solution}

\begin{solution}{exo:b3:extension-strike-slip:5}
$563/20 = 28$ tot $563/15 = \qty{38}{mm/yr}$. De verplaatsing is een minimumwaarde, en
de drift die door landmeting over de breukzone is gemeten, omvat ook beweging op
naburige breuken van het systeem en elastische rek die tussen
aardbevingen is opgeslagen.
\end{solution}

\begin{solution}{exo:b3:extension-strike-slip:6}
$t = 90 - \text{helling}$. Bij \ang{60}: $(1 + 0.6\times1.732)/(1 - 0.6\times
0.577) = 2.039/0.654 = 3.12$; bij \ang{50}: $1.715/0.497 = 3.45$; bij
\ang{40}: $1.504/0.285 = 5.28$. Vastlopen bij $\arctan0.6 = \ang{31.0}$.
\end{solution}

\begin{solution}{exo:b3:extension-strike-slip:7}
$\arctan0.4 = \ang{21.8}$ en $\arctan0.2 = \ang{11.3}$. Een breuk met een helling van
\ang{15} schuift alleen als $\mu < \tan\ang{15} = 0.27$.
\end{solution}

\begin{solution}{exo:b3:extension-strike-slip:8}
$2.5\times4.0 = \qty{10.0}{km^2}$ en $2.5\times(4.0 - 1.5) =
\qty{6.25}{km^2}$.
\end{solution}

\begin{solution}{exo:b3:extension-strike-slip:9}
$132/1.19 = \qty{110.9}{km}$, uitgerekt met \qty{21.1}{km}, wat
$21.1/6 = \qty{3.5}{Myr}$ kost.
\end{solution}

\begin{solution}{exo:b3:extension-strike-slip:10}
$3.0/\sin\ang{60} = \qty{3.464}{km}$ vóór en $3.0/\sin\ang{45} =
\qty{4.243}{km}$ na: \qty{0.779}{km} per blok ($\beta = 1.22$). De lagen
hellen $60 - 45 = \ang{15}$. Elk blok draait als geheel: de rotatie die
de helling van breuken die, bijvoorbeeld, naar het oosten hellen, verlaagt, draait een horizontale lijn in het
blok zo dat die naar het oosten stijgt, en de lagen hellen naar het westen, tegengesteld aan de
breuken.
\end{solution}

\begin{solution}{exo:b3:extension-strike-slip:11}
In een sinistrale oost-westzone loopt $\sigma_1$ op \ang{045}. R-schuifbreuken op
$\phi/2 = \ang{20}$ van de zone aan de kant van $\sigma_1$: strekking \ang{070},
sinistraal. R$'$ op $90 - 20 = \ang{70}$: strekking \ang{020}, dextraal. Trekscheuren
evenwijdig aan $\sigma_1$: strekking \ang{045}.
\end{solution}

\begin{solution}{exo:b3:extension-strike-slip:12}
Voor $\mu = 0.6$ ontstaan de breuken op \ang{60.5} en lopen ze vast op \ang{31.0}: $\beta =
\sin\ang{60.5}/\sin\ang{31.0} = 0.870/0.515 = 1.69$. Voor $\mu = 0.85$: van
\ang{65.2} tot \ang{40.4}, $\beta = 0.908/0.648 = 1.40$. Breuken met een lagere
wrijving roteren verder voordat ze vastlopen: zwakke gesteenten rekken meer uit voordat
er nieuwe breuken moeten ontstaan.
\end{solution}

\begin{solution}{pb:b3:extension-strike-slip:1}
\textbf{1.} De lagen waren horizontaal en zijn met de blokken meegedraaid; de
breuken zijn over dezelfde hoek gedraaid, dus $\theta_0 = 42 + 18 = \ang{60}$,
Andersons helling voor $\phi = \ang{30}$.

\textbf{2.} $\beta = \sin\ang{60}/\sin\ang{42} = 0.8660/0.6691 = 1.294$.

\textbf{3.} $24.0/1.294 = \qty{18.5}{km}$; het domein is met
\qty{5.5}{km} uitgerekt (\qty{5.46}{km}).

\textbf{4.} $t = \ang{48}$: $(1 + 0.6\times0.900)/(1 - 0.6\times1.111) =
1.540/0.334 = 4.62$, tegen $3.12$ voor nieuwe breuken op \ang{60.5}: ruim
de helft meer.

\textbf{5.} Bij $\arctan0.6 = \ang{31.0}$; $\beta = \sin\ang{60}/\sin\ang{31}
= 1.68$ hooguit. Daarvoor zouden nieuwe breuken op ongeveer \ang{60} de
blokken doorsnijden, en zou een tweede generatie dominostructuren het overnemen.

\textbf{6.} Elke kolom van de hangende wand daalt met de val van de breuk
eronder over de horizontale verplaatsing; de breuk is het steilst nabij het oppervlak en wordt er
vlakker vanaf, dus de kolommen dicht bij de breuk dalen het meest en die ver weg
het minst: de lagen kantelen naar de breuk toe.

\textbf{7.} $D = A/H = 16.0/3.2 = \qty{5.0}{km}$.

\textbf{8.} $D$ is evenredig met $A$: $\qty{5.0}{km} \pm \qty{0.5}{km}$.

\textbf{9.} $3.2\times(5.0 - 2.0) = \qty{9.6}{km^2}$.

\textbf{10.} De horizontale verplaatsing, \qty{3.2}{km}.

\textbf{11.} De vulling wordt dikker naar de breuk toe, in een wig, en de oudere
lagen zijn sterker gekanteld dan de jongere: elke laag werd afgezet terwijl
de rollover en de ruimte naast de breuk bleven groeien.

\textbf{12.} $z = P/(\rho g) = \num{5.5e8}/(2800\times9.8) = \qty{20.0}{km}$.

\textbf{13.} $20.0/(24 - 16) = \qty{2.5}{km/Myr} = \qty{2.5}{mm/yr}$.

\textbf{14.} $20.0/\sin\ang{30} = \qty{40}{km}$ aan verschuiving, en $40\cos\ang{30} =
\qty{34.6}{km}$ aan horizontale rek.

\textbf{15.} Nee: \ang{15} ligt onder de vergrendelingshelling van \ang{31}. Ofwel gleed de
breuk toen ze steiler was en werd ze daarna gekanteld (\omterm{def:b3:extension-strike-slip:core-complex}{rollend scharnier}), ofwel
was ze zwak, met $\mu < 0.27$ door klei of vloeistoffen onder hoge druk.

\textbf{16.} Door $45 - 15 = \ang{30}$: terwijl de bovenplaat wegbewoog, rees de
ontlaste liggende wand door isostasie op en boog door, waardoor de breuk via het
scharnier werd gekanteld.

\textbf{17.} \qty{34.6}{km} tegenover de \qty{8.7}{km} van de halve slenk (vraag
23): viermaal zoveel. Kerncomplexen liggen waar de rek het grootst is, zoals in
het centrale Basin and Range.

\textbf{18.} Een linkse horizontale verschuiving die naar links springt: extensief, er ontstaat een bekken.

\textbf{19.} $\qty{5}{mm/yr}\times\qty{1.2}{Myr} = \qty{6.0}{km}$.

\textbf{20.} $4.0\times6.0 = \qty{24}{km^2}$; $3.0 + 6.0 = \qty{9.0}{km}$ lang;
$9.0/4.0 = 2.25$.

\textbf{21.} $132/10 = 13$ tot $132/7 = 19$: het Dode Zee-bekken is veel
langgerekter, de verplaatsing is vele malen de breedte.

\textbf{22.} Een \omterm{def:b3:extension-strike-slip:bends}{compressieve bocht}: een rug die tussen de takken wordt opgeduwd; in
doorsnede een positieve \omterm{def:b3:extension-strike-slip:flower}{bloemstructuur}, waarbij de breuk naar boven uitwaaiert in takken
met een overschuivende component, zoals de takken van 100 tot \qty{300}{m} breed op de
Dode Zee Transform.

\textbf{23.} $3.2 + 5.46 = \qty{8.66}{km}$: de halve slenk van Merrow is
verlengd met \textbf{\qty[text-series-to-math=true]{8.7}{km}}.
\end{solution}
